\section{Detector and Evaluation Setup}\label{sec:setup}

\emph{Data.} We use ten CIC-IoT2023 captures~\cite{neto2023ciciot2023}: Benign, DDoS-SYN, DoS-HTTP, Recon-PortScan, MITM-ArpSpoofing, DictionaryBruteForce, SqlInjection, XSS, CommandInjection and Mirai-udpplain (about 5.17\,GB). Files over 300\,MB are trimmed to their first $10^6$ packets. The captures map to eight classes: Benign, BruteForce, DDoS-SYN, DoS-HTTP, MITM-ARP, Mirai, Recon and Web (SQLi, XSS and command injection).

\emph{Flows and features.} Flows are built with NFStream~6.6.0~\cite{aouini2022nfstream} (\texttt{statistical\_analysis=True}, idle timeout 15\,s, active timeout 120\,s), with identical settings offline and on the live gateway. The model uses 54 statistical features (packet counts and sizes, inter-arrival times, durations, TCP flag counts per direction). Identifiers (IP and MAC addresses, ports, protocol, timestamps) are excluded, because they let a model memorise the capture~\cite{arp2022dos}. This yields 857,206 labelled flows, 511,527 of them DDoS-SYN (a consequence of trimming).

\emph{Detector.} The detector is a scikit-learn~\cite{pedregosa2011sklearn} decision tree with balanced class weights; we compare it with a random forest (RF, 100 trees) and histogram gradient boosting (HGB) on the same features. We report two trees: the unconstrained baseline tree (temporal macro-F1 0.731, Benign FPR 3.9\%) in Section~\ref{sec:splits}, and the deployed tree, with limits tuned on validation (0.725, Benign FPR 3.8\%), everywhere else.

\emph{Splits.} The \emph{temporal} split puts the first 70\% of each capture's flows, by start time, into training (600,050 training and 257,156 test flows). The \emph{stratified random} split draws 70\% at random, stratified by class. Every design choice (tree limits, feature subset, blocking policy) is made on a validation set, the last 20\% of each capture's training period; the test period is used only for reporting.

\emph{Metrics.} We report per-class F1 with test support, macro-F1 over the eight classes, and the Benign false-positive rate (FPR), the share of Benign test flows predicted as any attack. Flows from one host pair are not independent, so 95\% confidence intervals come from a \emph{host-pair bootstrap}: 1,000 resamples of whole (capture, host pair) clusters (3,465 clusters on the temporal and 3,858 on the random test set). A flow-level bootstrap~\cite{efron1994bootstrap} gives far too narrow intervals: [0.712, 0.742] for the temporal decision tree, against [0.577, 0.823] with host pairs.

\emph{Model and size choices.} The three model families are close on the temporal split: tree 0.731, RF 0.737 at about 94$\times$ the size and 58$\times$ the single-flow latency (20,849 vs 222\,KiB; 7.826 vs 0.136\,ms per call), and HGB 0.714 with a Benign FPR of 7.6\%. We therefore deploy one tree. Limiting depth and leaf size on validation shrank it but raised the validation Benign FPR (from 1.3\% to 5.7\% at depth 16, leaf 10), so the chosen tree (depth 32, leaf 1) is effectively unconstrained: temporal macro-F1 0.725, 1,767 nodes, 222\,KiB. NSGA-II feature selection~\cite{deb2002nsga2,blank2020pymoo} picked 14 of 54 features on validation at no measurable cost (0.723) but barely shrank the tree (1,681 nodes, 211\,KiB): it reduces inputs, not model size.
